%% ============================================================
%%  §11  Related Work
%% ============================================================

\section{Related Work}
\label{sec:related}

\subsection{Rust Verification in Lean 4}

\textbf{Aeneas}~\cite{ho2022aeneas} is the primary related system. It translates
safe Rust programs via Charon into pure Lean~4 (and F$^\star$, Coq, HOL4)
using a shallow monadic embedding. Aeneas and lean-pl-verify share the goal
of Lean~4-based Rust verification and the use of a state-error monad, but
they differ fundamentally in embedding strategy. In Aeneas, the Lean function
generated for a Rust function \emph{is} the functional model; its correctness
depends on the soundness of the Charon+Aeneas pipeline. In lean-pl-verify,
the Lean definition of \texttt{evalStmtFuel} \emph{is} the semantics; there
is nothing to be sound with respect to, beyond what the kernel computes.

Recent applications of Aeneas include verified cryptographic
primitives~\cite{aeneascrypto2026}: the BLAKE3, SHA-2, and Kyber
implementations in the HACL* library. These verify real-world, production
code at a scale that lean-pl-verify does not currently match. The advantage
of Aeneas in this context is the Charon pipeline's ability to extract from
arbitrary crates; our manual transcription limits scale.

\textbf{LOOM}~\cite{loom2026} develops a framework for multi-modal
foundational verification using Dijkstra monads in Lean~4. LOOM's goals
overlap with ours (language-agnostic specifications, Lean~4 kernel as trust
anchor), but LOOM operates on shallow program representations rather than
deep AST embeddings. The Dijkstra monad index provides automated WP
generation, which LOOM exploits for more automated proofs. Whether LOOM's
approach can be extended to cross-language unification theorems similar to
ours is an interesting open question.

\subsection{Rust Verification via SMT}

\textbf{Verus}~\cite{verus2023} is currently the most practical tool for
verifying real-world Rust programs, with successful applications to distributed
systems and operating system code. Verus embeds specifications directly into
Rust source code using a ghost-type discipline; the Z3 SMT solver handles
proof obligations. The principal limitation is the ``SMT cliff'': complex
proofs, especially those involving non-linear arithmetic or intricate data
structure invariants, can cause Z3 to time out. Verus does not support
TypeScript and does not produce Lean kernel-checked proofs.

\textbf{Creusot}~\cite{creusot2022} uses Why3 as a deductive verification
backend. Programs are annotated with Rust attributes, and Creusot generates
Why3 goals discharged by a portfolio of SMT solvers. It supports a wider
fragment of Rust than Verus (notably, trait objects and iterators), but
shares the reliance on SMT solvers in the TCB.

\subsection{Bounded Verification}

\textbf{Kani}~\cite{kani2022} is a bounded model checker for Rust built on
CBMC. It handles unsafe Rust and concurrency by bounding the execution depth
and encoding safety properties as assertions for a SAT/SMT solver. Kani
excels at finding concrete bugs (memory unsafety, integer overflow in specific
inputs) but cannot prove properties for all inputs or for loops with
unbounded iteration counts. It does not produce formal proof certificates.

\textbf{SEABMC} and related tools apply bounded symbolic execution to
Rust via LLVM IR. They share Kani's strengths and limitations regarding
unbounded verification.

\subsection{Separation Logic and Unsafe Rust}

\textbf{RustBelt}~\cite{rustbelt2018} is the foundational semantic model
for Rust's type system, developed in Coq using the Iris separation logic
framework~\cite{iris2015}. RustBelt proves that the Rust type system
enforces memory safety even in the presence of unsafe code, by providing
a semantic model for lifetimes and type invariants. The proofs are highly
manual and cover only a small fragment of the standard library.

Our work shares RustBelt's interest in formal semantics for Rust but
targets a complementary level: functional correctness of safe programs
rather than type-soundness of unsafe abstractions. Combining a
RustBelt-style unsafe-code foundation with our safe-code verification
layer would be a natural long-term goal, following the hybrid architecture
proposed by Jung~\etal~\cite{rustbelt2018}.

\subsection{Cross-Language and Multi-Language Verification}

Cross-language correctness has been studied in the context of compilation
verification~\cite{compcert2009}: if a compiler is verified to preserve
semantics, then properties proved about the source program hold for the
target. Our approach is different: we do not verify a compiler; we instead
provide a common semantic target into which multiple source languages are
independently embedded.

\textbf{Bedrock2} and related work in the certified programming languages
community has explored ``push-button'' verification of C code compiled to
machine code, establishing end-to-end correctness from source to binary.
The multi-language direction in lean-pl-verify is orthogonal: rather than
going deeper (closer to hardware), we go wider (across languages).

\textbf{F$^\star$}~\cite{fstar2016} supports multiple languages via a
common F$^\star$ extraction target. The Low$^\star$ sublanguage is used
to verify C programs; KReMLin extracts to C. This is structurally similar
to our approach (common monad, multiple source embeddings) but F$^\star$'s
multi-language support is directed at compilation rather than cross-language
property sharing.

\subsection{Proof Automation for Program Verification}

Lean~4's \texttt{grind} tactic~\cite{lean4} combines congruence closure,
E-matching, and linear arithmetic; it is designed precisely for program
verification goals. We deliberately avoided \texttt{grind} in lean-pl-verify,
preferring explicit \texttt{simp only} calls that make the proof state
predictable. Whether \texttt{grind} could further compress the loop invariant
proofs is unexplored.

The Lean Squad project~\cite{leansquad2026} reports near-zero human labour
verification of Rust programs using an LLM agent that drives the Lean LSP
via the Model Context Protocol. This is architecturally different from our
work (agent vs.\ manual proof writing) but shares the goal of using Lean~4
as the verification target. The agent uses Aeneas for the Rust-to-Lean
translation; combining the agent with our deep embedding would be an
interesting experiment.

\subsection{Domain-Specific Languages for Verification}

\textbf{Dafny}~\cite{dafny2010} provides an integrated verification language
where specifications (\texttt{ensures}/\texttt{requires}/\texttt{invariant})
are interleaved with program code. It supports loops, recursive functions,
and SMT-backed automation. Dafny does not target Rust or TypeScript.

\textbf{DSLean}~\cite{dslean2026} develops a framework for type-correct
interoperability between Lean~4 and external DSLs, using Lean's elaboration
and metaprogramming to enforce typing in real time. This is related to our
macro-based proof automation (\S\ref{sec:automation}) but aimed at DSL
interoperability rather than proof compression.
